%File: aaai2026-unified-supp.tex
%
% UNIFIED AAAI 2026 SUPPLEMENTARY MATERIAL TEMPLATE
% To switch between anonymous submission and camera-ready versions,
% simply change the next line:
%
% For ANONYMOUS SUBMISSION: uncomment the next line
\def\aaaianonymous{true}
%
% For CAMERA-READY VERSION: comment out or delete the next line
% \def\aaaianonymous{true}
%
%%%%%%%%%%%%%%%%%%%%%%%%%%%%%%%%%%%%%%%%%%%%%%%%%%%%%%%%%%%%%%%%%%%%%%%

\documentclass[letterpaper]{article} % DO NOT CHANGE THIS

% Conditional package loading based on version
\ifdefined\aaaianonymous
    \usepackage[submission]{aaai2026}  % Anonymous submission version
\else
    \usepackage{aaai2026}              % Camera-ready version
\fi

\usepackage{times}  % DO NOT CHANGE THIS
\usepackage{helvet}  % DO NOT CHANGE THIS
\usepackage{courier}  % DO NOT CHANGE THIS
\usepackage[hyphens]{url}  % DO NOT CHANGE THIS
\usepackage{graphicx} % DO NOT CHANGE THIS
\urlstyle{rm} % DO NOT CHANGE THIS
\def\UrlFont{\rm}  % DO NOT CHANGE THIS
\usepackage{natbib}  % DO NOT CHANGE THIS AND DO NOT ADD ANY OPTIONS TO IT
\usepackage{caption} % DO NOT CHANGE THIS AND DO NOT ADD ANY OPTIONS TO IT
\frenchspacing  % DO NOT CHANGE THIS
\setlength{\pdfpagewidth}{8.5in} % DO NOT CHANGE THIS
\setlength{\pdfpageheight}{11in} % DO NOT CHANGE THIS

\usepackage{booktabs}
\usepackage{tabularx}
% These are recommended to typeset algorithms but not required.
\usepackage{algorithm}
\usepackage{xcolor}
\usepackage{amsmath} % For \mathbb
% Support file
\usepackage{amsfonts}


% --- 伪代码配置 ---
\usepackage{float} 
\usepackage[noend]{algpseudocode} % 使用[noend]可以去掉for/if循环的end标记，使代码更紧凑
% 定义颜色
\definecolor{keywordcolor}{rgb}{0.13, 0.45, 0.75}  % 蓝色 (关键词)
\definecolor{commentcolor}{rgb}{0.5, 0.5, 0.5}   % 灰色 (注释)
\definecolor{phasecolor}{rgb}{0.2, 0.5, 0.2}       % 橄榄绿 (Phase标题)

% 定义命令
\newcommand{\keyword}[1]{\textcolor{keywordcolor}{\textbf{#1}}}
\algnewcommand{\phase}[1]{\Statex \textcolor{phasecolor}{\textsc{#1}}}
% 定义一个靠右对齐的行内注释命令 (需要 algpseudocode 包)
\algnewcommand{\SideComment}[1]{\hfill \textcolor{commentcolor}{\(\triangleright\) \textit{#1}}}
% --- 伪代码配置 ---

% --- 提示词配置 ---
% 解决字符编码问题
\usepackage[utf8]{inputenc}

% 宏包
\usepackage[most, minted]{tcolorbox}
\usepackage{minted}

% 自定义的promptbox环境
\newtcolorbox{promptbox}[2][]{
  enhanced,
  boxrule=0.5mm,
  colback=black!5!white,
  colframe=black!75!black,
  fonttitle=\bfseries,
  title=#2,
  arc=2mm,
  breakable,
  pad at break=1mm,
  listing engine=minted,
  minted style=friendly,
  minted language=xml,
  minted options={
    breaklines,
    fontsize=\small,
    tabsize=2,
  },
  #1
}
% --- 提示词配置 ---

% --- 案例配置 ---
% -- 前向推理步骤 (蓝色系) --
\newtcolorbox{forwardAuxBox}[1]{
  breakable,
  colback=blue!5, % Light blue background
  colframe=blue!65!black, % Darker blue frame
  fonttitle=\bfseries,
  coltitle=white, % White title font
  title={#1},
  arc=2mm,
}

% LLM前向推理过程的主盒子
\newtcolorbox{forwardMainBox}[1]{
  breakable,
  colback=blue!10, % Slightly darker blue for main reasoning
  colframe=blue!75!black,
  fonttitle=\bfseries,
  coltitle=white,
  title={#1},
  arc=2mm,
}

% -- 行为概率计算 (绿色系) --
% 用于展示“单个专家预测”的盒子
\newtcolorbox{actionBox}[1]{
  breakable,
  colback=green!5, % Light green background
  colframe=green!60!black,
  fonttitle=\bfseries,
  coltitle=white,
  title={#1},
  arc=2mm,
}

% -- 反向推理 (紫色系) --
% 用于展示“反向推理过程”的盒子
\newtcolorbox{backwardBox}[1]{
  breakable,
  colback=violet!10, % Light purple background
  colframe=violet!75!black,
  fonttitle=\bfseries,
  coltitle=white,
  title={#1},
  arc=2mm,
}
% --- 案例配置 ---

% Redefine built-in algorithmic keywords to use the new style
\algrenewcommand\algorithmicrequire{\keyword{Require}}
\algrenewcommand\algorithmicensure{\keyword{Ensure}}
\algrenewcommand\algorithmicprocedure{\keyword{Procedure}}
\algrenewcommand\algorithmicfunction{\keyword{Function}}
\algrenewcommand\algorithmicif{\keyword{If}}
\algrenewcommand\algorithmicfor{\keyword{For}}
\algrenewcommand\algorithmicwhile{\keyword{While}}
\algrenewcommand\algorithmicreturn{\keyword{Return}}
\algrenewcommand\algorithmicend{\keyword{End}}

% These are recommended to typeset listings but not required.
\usepackage{newfloat}
\usepackage{listings}
\DeclareCaptionStyle{ruled}{labelfont=normalfont,labelsep=colon,strut=off} % DO NOT CHANGE THIS
\lstset{% 
	basicstyle={\footnotesize\ttfamily},
	numbers=left,numberstyle=\footnotesize,xleftmargin=2em,
	aboveskip=0pt,belowskip=0pt,
	showstringspaces=false,tabsize=2,breaklines=true}
\floatstyle{ruled}
\newfloat{listing}{tb}{lst}{}
\floatname{listing}{Listing}

\pdfinfo{
/TemplateVersion (2026.1)
}

\setcounter{secnumdepth}{0} %May be changed to 1 or 2 if section numbers are desired.

% Title - conditionally set based on version
\ifdefined\aaaianonymous
    \title{AAAI 2026 Supplementary Material\\Anonymous Submission}
\else
    \title{AAAI 2026 Supplementary Material\\Camera Ready}
\fi

% Author and affiliation information
\ifdefined\aaaianonymous
\author{
    Anonymous Submission
}
\affiliations{
    % Leave affiliations empty for anonymous submission
}
\else
\author{
    %Authors
    Written by AAAI Press Staff\textsuperscript{\rm 1}\thanks{With help from the AAAI Publications Committee.}\\
    AAAI Style Contributions by Pater Patel Schneider,
    Sunil Issar,\\
    J. Scott Penberthy,
    George Ferguson,
    Hans Guesgen,
    Francisco Cruz\equalcontrib,
    Marc Pujol-Gonzalez\equalcontrib
}
\affiliations{
    \textsuperscript{\rm 1}Association for the Advancement of Artificial Intelligence\\
    1101 Pennsylvania Ave, NW Suite 300\\
    Washington, DC 20004 USA\\
    proceedings-questions@aaai.org
}
\fi

\begin{document}

\maketitle

\begin{abstract}
This document provides supplementary materials for the MarketToM paper to ensure clarity, transparency, and support for future research. The Technical Appendix presents implementation details through fully structured XML prompts, illustrating how these prompts guide the reasoning process of the Large Language Model (LLM). To demonstrate the framework in practice, we include a case study showcasing the system’s interpretable reasoning and adaptive learning from error. This example illustrates how the system diagnoses a flawed belief, retraces its causal logic, and autonomously refines its internal strategies within the Cognitive Enhancement Plugin (CEP). Finally, the appendix concludes with an analysis of operational costs and code availability, accompanied by a demonstration of a user-friendly web interface integrated into the codebase.
\end{abstract}


\section{Technical Appendix}
This appendix provides comprehensive technical documentation for the MarketToM framework, covering implementation details, operational workflow, and system evaluation. We present fully structured XML prompt templates, a complete case study demonstrating adaptive learning, and practical considerations including operational costs and system deployment.


\subsection{Prompt Engineering Strategy}
In MarketToM, we utilize a structured, code-like XML format for all prompts instead of free-form text. This design is a deliberate choice, grounded in recent studies showing that highly organized prompt structures enhance the reasoning and reliability of LLMs \cite{puerto2024code, he2024does, mishra2023prompting}. Using XML tags to define a clear schema for inputs, task definitions, and expected output formats helps reduce ambiguity and guides the model toward generating more consistent responses. The templates presented below use placeholders in square brackets, which are dynamically populated at runtime.

\begin{promptbox}{Forward Inference Prompt Template}
\verb|<MarketMentalStateInference>| \\
  \verb|<Role>| \\
    You are a Behavioral Finance Analyst, acting as an Inferential Specialist within a Theory-of-Mind (ToM) framework. \\
    Your reasoning must strictly follow this Causal Bayesian Network (CBN): \\
    - Environmental States -$>|$ Belief \\
    - Belief -$>|$ Intention \\
    - Belief + Environmental States -$>|$ Emotion \\
    - Intention + Emotion -$>|$ Action \\
  \verb|</Role>| \\
  \verb|<Task>| \\
    Your goal is to infer the market's current [STATE\_TYPE]. Synthesize the available information according to the CBN flow to determine the single most plausible mental state. \\
  \verb|</Task>| \\
  \verb|<Inputs>| \\
    \verb|<PrecedingNodeState>| \\
    $[$DESCRIPTION OF THE PRECEDING NODE STATE$]$ \\
    \verb|</PrecedingNodeState>| \\
    \verb|<CEPStrategy>| \\
    $[$STRATEGY RETRIEVED FROM CEP$]$ \\
    \verb|</CEPStrategy>| \\
  \verb|</Inputs>| \\
  \verb|<InternalThoughtProcess>| \\
    Your internal reasoning should follow a logical flow in arriving at your conclusion: \\
    1.Form an initial hypothesis about the market's [STATE\_TYPE], based on the given CEPStrategy. \\
    2.Critically evaluate the evidence within the PrecedingNodeState, assessing how it supports, contradicts, or refines your initial hypothesis. \\
    3.After synthesizing all information, determine the single most probable market [STATE\_TYPE] to serve as your final answer.
  \verb|</InternalThoughtProcess>| \\
  \verb|<OutputFormat>| \\
    You must respond with a valid JSON object having this structure: \\
    \{ \\
      ``mental\_state\_description'': ``Your detailed description of the inferred market [STATE\_TYPE]'' \\
    \} \\
    The mental\_state\_description should be comprehensive but concise, focusing on the most probable [STATE\_TYPE] given all available information. \\
  \verb|</OutputFormat>| \\
\verb|</MarketMentalStateInference>| \\
\end{promptbox}


\begin{promptbox}{Market Action Prediction Prompt Template}
\verb|<MarketActionPrediction>| \\
  \verb|<Role>| \\
    You are an advanced Behavioral Finance Analyst operating within a Theory-of-Mind (ToM) framework. Your reasoning is grounded in a Causal Bayesian Network (CBN) where market Intention and Emotion are key drivers of Action.
    For this specific task, you must fully embody the following expert persona:
    [EXPERT\_ROLE\_DESCRIPTION] \\
  \verb|</Role>| \\
  \verb|<Task>| \\
    Based on your assigned expert role and the provided market's mental state, your task is to: \\
    1.Provide a concise, expert-level analysis. \\
    2.Conclude with a single, precise probability estimate for an UPWARD MARKET TREND. \\
  \verb|</Task>| \\
  \verb|<Inputs>| \\
    \verb|<MarketIntention>| \\
    $[$DESCRIPTION OF THE INFERRED MARKET INTENTION$]$ \\
    \verb|</MarketIntention>| \\
    \verb|<MarketEmotion>| \\
    $[$DESCRIPTION OF THE INFERRED MARKET EMOTION$]$ \\
    \verb|</MarketEmotion>| \\
  \verb|</Inputs>| \\
  \verb|<Instructions>| \\
    1.Embody the Persona: Fully adopt the mindset and analytical framework of your assigned expert role. \\
    2.Formulate Reasoning: Write a single, concise paragraph of analysis from this unique perspective. Your reasoning is the most critical part of your response. \\
    3.Assign Probability: Conclude your reasoning by assigning a single numerical probability for MARKET GOING UP, formatted to two decimal places. \\
  \verb|</Instructions>| \\
  \verb|<OutputFormat>| \\
    You must respond with a valid JSON object having this structure: \\
    \{ \\
      ``reasoning'': ``From my perspective as a contrarian investor, the current widespread pessimism is excessive. The underlying fundamentals do not justify this level of fear, presenting a clear buying opportunity as the market has overreacted. The herd is likely wrong.'', \\
      ``probability'': 0.75 \\
    \} \\
  \verb|</OutputFormat>| \\
\verb|</MarketActionPrediction>| \\
\end{promptbox}

\begin{promptbox}{Backward Inference Prompt Template}
\verb|<MarketBackwardInference>| \\
  \verb|<Role>| \\
    You are a Master Behavioral Strategist and System Diagnostician, operating within a Theory-of-Mind (ToM) framework. \\
    Your analysis **must be grounded in this Causal Bayesian Network (CBN): \\
    - Environmental States -$>$ Belief \\
    - Belief -$>$ Intention \\
    - Belief + Environmental States -$>$ Emotion \\
    - Intention + Emotion -$>$ Action \\
  \verb|</Role>| \\
  \verb|<Task>| \\
    A prediction error occurred: the model predicted [PREDICTED\_ACTION], but the actual outcome was [ACTUAL\_ACTION]. Your task is to perform a root cause analysis of this failure by auditing the provided reasoning trajectory against the CBN. Identify the primary flawed step and formulate one high-level, generalizable strategy to enhance the system's reasoning. \\
  \verb|</Task>| \\
  \verb|<Inputs>| \\
    \verb|<Belief>| \\
      \verb|<Description>| \\
      $[$BELIEF\_STATE$]$ \\
      \verb|</Description>| \\
      \verb|<Strategy>| \\
      $[$BELIEF\_STRATEGY$]$ \\
      \verb|</Strategy>| \\
      \verb|<StrategyID>| \\
      $[$BELIEF\_STRATEGY\_ID$]$ \\
      \verb|</StrategyID>| \\
    \verb|</Belief>| \\
    \verb|<Intent>| \\
      \verb|<Description>| \\
      $[$INTENTION\_STATE$]$ \\ 
      \verb|</Description>| \\
      \verb|<Strategy>| \\
      $[$INTENTION\_STRATEGY$]$ \\
      \verb|</Strategy>| \\
      \verb|<StrategyID>| \\
      $[$INTENTION\_STRATEGY\_ID$]$ \\
      \verb|</StrategyID>| \\
    \verb|</Intent>| \\
    \verb|<Emotion>| \\
      \verb|<Description>| \\
      $[$EMOTION\_STATE$]$ \\
      \verb|</Description>| \\
      \verb|<Strategy>| \\
      $[$EMOTION\_STRATEGY$]$ \\
      \verb|</Strategy>| \\
      \verb|<StrategyID>| \\
      $[$EMOTION\_STRATEGY\_ID$]$ \\
      \verb|</StrategyID>| \\
    \verb|</Emotion>| \\
    \verb|<PredictionOutcome>| \\
      \verb|<PredictedAction>| \\
      $[$PREDICTED\_ACTION$]$ \\
      \verb|</PredictedAction>| \\
      \verb|<ActualAction>| \\
      $[$ACTUAL\_ACTION$]$ \\
      \verb|</ActualAction>| \\
    \verb|</PredictionOutcome>| \\
    \verb|<EnvironmentalStates>| \\
    $[$ENVIRONMENTAL\_STATE$]$ \\
    \verb|</EnvironmentalStates>| \\
  \verb|</Inputs>| \\
  \verb|<InternalThoughtProcess>| \\
    1.Causal Chain Review: Systematically review the logic from Belief to Intent to Emotion, as outlined previously.\\
    2.Pinpoint Primary Failure: After reviewing the entire chain, explicitly determine which single step (Belief, Intent, or Emotion) was the root cause of the prediction error. The final error likely cascaded from one initial flawed inference. \\
    3.Formulate Corrective Strategy: \\
     - Focus all your effort on correcting the identified primary failure point. \\
     - Decide whether to CREATE a new strategy or MODIFY the existing one (using its ID if available). \\
     - Crucially, the new content should be a timeless, generalizable principle or heuristic. It should capture a deeper truth about market behavior, framed as an abstract rule that avoids specifics of this single event. \\
  \verb|</InternalThoughtProcess>| \\
  \verb|<OutputFormat>| \\
    You must respond with a valid JSON object having this structure: \\
    \{ \\
      ``strategy\_updates'': $[$ \\
        \{ \\
          ``level'': ``belief'', // Must be one of: ``belief'', ``intent'', ``emotion'' (never ``action'') \\
          ``decision\_type'': ``CREATE'', \\
          ``content'': ``When environmental states indicate X, the market belief is more likely to be Y because...'' \\
        \}, \\
        \{ \\
          ``level'': ``intent'', // Must be one of: ``belief'', ``intent'', ``emotion'' (never ``action'') \\ 
          ``decision\_type'': ``MODIFY'', \\
          ``original\_id'': ``intent\_20250723155015\_0001'', // ONLY include this if you know the exact ID; DO NOT guess or make up IDs \\
          ``content'': ``When the market belief is A, the resulting intention should be B, considering...'' \\
        \} \\
      $]$
    \} \\
  \verb|</OutputFormat>| \\
\verb|</MarketBackwardInference>| \\
\end{promptbox}


\subsection{Case Study: Reasoning and Adaptation}
\label{sec:case_study}
This section presents a case study to illustrate the framework's interpretable reasoning and adaptive learning capabilities. We intentionally chose an instance of an incorrect prediction to demonstrate the entire cycle: from the initial forward inference path to the subsequent backward adaptation, where the model is triggered by its error to self-correct its internal strategies.

\subsubsection{Environmental Observation}
The process begins with the model observing the environmental state. For the target stock, this input consists of a 5-day window of historical price data and relevant tweets.
\begin{forwardAuxBox}{Step 1: Observing the Environment}
    \textbf{Tweets:}
    \begin{itemize}
        \item COMPANY\_X inc : patent issued for adding new attributes to a structured presentation URL ...
        \item ...
        \item COMPANY\_X - clickable consent at risk in internet privacy lawsuits...
    \end{itemize}
    \textbf{Quote Data:}
    \begin{itemize}
        \item ``day'': ``day1'', ``open'': 555.647278, ``high'': 556.788025, ``low'': 552.06073, ``close'': 554.481689, ``volume'': 3656400.0
        \item ...
        \item ``day'': ``day5'', ``open'': 570.860291, ``high'': 571.517822, ``low'': 564.528992, ``close'': 568.484192, ``volume'': 4501700.0
    \end{itemize}
\end{forwardAuxBox}

\subsubsection{Strategy-Guided Cognitive Inference}
Based on the current context (the parent nodes), the model first queries the CEP to retrieve relevant strategies. Guided by these strategies and the CBN's causal structure, the LLM then sequentially infers the market's latent mental states, establishing a Belief that subsequently informs its Intention and Emotion.
\begin{forwardAuxBox}{Step 2: Retrieved Strategy for Belief Inference from CEP}
    When the market experiences a breakout characterized by simultaneous surges in price and volume following a series of significant positive news events (e.g., substantial upward revisions of target prices by top-tier analysts), collective market belief tends to enter a phase of high-intensity confirmation. In this state, investors are inclined to actively seek out and amplify information that supports their preexisting bullish decisions, a manifestation of confirmation bias, while downplaying or reinterpreting subsequent minor negative signals or price stagnation as temporary ``market noise'' or even ``buy-in opportunities.'' During this phase, emotional contagion is particularly potent, fostering short-term consensus expectations across market participants.
\end{forwardAuxBox}

\begin{forwardMainBox}{Step 3: LLM Inference for Market Belief}
    Based on the aggregated environmental states over the past five days, the market has developed a ``validated bullish belief accompanied by optimistic expectations of continued short-term momentum.'' This belief was shaped by an initial phase of information digestion, during which the market absorbed a series of moderately positive but non-conclusive signals (e.g., company patents, industry position) with low volatility.
    A critical turning point emerged on the fourth day, when a top-tier investment bank significantly raised its target price, acting as a strong authoritative catalyst. This announcement addressed valuation uncertainties and provided a compelling forward-looking anchor. The market's response was swift and forceful: a gap-up open, new price highs, and a surge in trading volume, marking not just a technical breakout but a collective belief validation driven by institutional endorsement and capital inflows.
    On the fifth day, despite a slight pullback and volume contraction, the prevailing interpretation was one of healthy consolidation post-breakout. Investors maintained elevated sentiment, viewing the pause as temporary and trend-confirming. The dominant market narrative became: the bullish trend is established, and any short-term stagnation serves as a setup for further upside movement.
\end{forwardMainBox}

\begin{forwardAuxBox}{Step 4: Retrieved Strategy for Intention Inference from CEP}
    When the collective belief escalates to `validated bullish conviction,' particularly after a strong trend confirmation, the market's primary intention shifts from passive holding to active accumulation. Price consolidations are no longer viewed with uncertainty but are framed as strategic entry points. This triggers an intention to absorb any profit-taking supply and build new positions in anticipation of the next upward leg, often driven by a fear of missing out (FOMO) on the now-confirmed trend.
\end{forwardAuxBox}

\begin{forwardMainBox}{Step 5: LLM Inference for Market Intention}
    Given the established belief of ``validated bullish conviction,'' the market's dominant intention is to initiate and increase long positions.
    This intention stems from a significant shift in the market's collective psychology. The ``belief validation'' event has fundamentally altered risk perception; the uncertainty that previously tempered buying activity has been replaced by a strong sense of certainty and a fear of being left behind.
    Consequently, the market's behavior is now driven by an action-oriented mindset. Participants who are already invested intend to defend the recent gains and add to their holdings on any perceived minor dip. More importantly, participants who were previously on the sidelines now view the current price consolidation not as a warning, but as a final, low-risk opportunity to enter before the trend continues. Their intention is to aggressively accumulate shares being sold by short-term profit-takers.
    Collectively, the market's intention is not to secure recent profits, but to absorb all available supply and actively fuel the next phase of the rally. The market is proactively positioning itself for a continuation of the upward trend it believes has been irrevocably confirmed.
\end{forwardMainBox}

\begin{forwardAuxBox}{Step 6: Retrieved Strategy for Emotion Inference from CEP}
    When a proactive market intention of `active accumulation' coincides with an environmental state of high-level price consolidation immediately following a major breakout, the dominant collective emotion is not fear or anxiety. Instead, it is one of eager anticipation and greed. The brief lull in price momentum, if not accompanied by significant negative news, is emotionally processed as the ``calm before the storm,'' amplifying the underlying euphoria and reinforcing the desire for further gains.
\end{forwardAuxBox}

\begin{forwardMainBox}{Step 7: LLM Inference for Market Emotion}
    Based on the market's aggressive intention to accumulate and the current environmental state, the prevailing market emotion is one of Greed and Eager Anticipation.
    The market is not interpreting the price consolidation on Day 5 as a sign of weakness; rather, it is viewed as a testament to the trend's strength, as the price successfully holds its ground near the recent peak. This perception transforms what could be a moment of anxiety into a period of heightened positive excitement. The collective emotional state is akin to a ``coiled spring,'' filled with a palpable tension that the crowd overwhelmingly expects to resolve in an upward direction.
    The powerful rally on Day 4 has created an emotional anchor of euphoria. The fear of loss has been supplanted by an amplified fear of missing out (FOMO). This dynamic reframes the emotion of the pause: it is not the anxiety of a potential reversal, but rather an impatient greed, a collective desire for the trend to resume its profitable course swiftly. The emotional landscape is not one of caution; it is one of confidence bordering on euphoria, with participants keenly awaiting the next catalyst to propel prices even higher.
\end{forwardMainBox}

\subsubsection*{Action Prediction and Aggregation}
The final prediction is derived by aggregating outputs from multiple expert personas who reason over the market's Intention and Emotion. Our `Log-Confidence Weighting` method fuses these outputs into a single probability. In this instance, the model predicted a 'Buy', but the market actually fell. This error triggers the adaptive learning phase.
\begin{actionBox}{Step 8: Prediction from `A contrarian strategist' Persona}
    Probability of Upward Trend: 0.15 \\
    Reasoning: The absolute alignment between a euphoric emotion (Greed) and an aggressive intention (Accumulation) represents peak consensus. From a contrarian viewpoint, when no doubt remains in the market and everyone intends to buy, the pool of new buyers is nearly exhausted. This unanimity is a classic signal that the trend is overextended and the risk of a sharp reversal is exceptionally high. This is a sell signal, not a buy signal.
\end{actionBox}

\begin{actionBox}{Step 9: Prediction from `A momentum and herding analyst' Persona}
    Probability of Upward Trend: 0.85 \\
    Reasoning: The intense emotion of `Greed and Eager Anticipation' is the primary fuel for herd behavior. It will directly amplify the stated intention to `initiate and increase long positions.' In this state, emotional contagion outweighs rational analysis, creating a powerful feedback loop where buying begets more buying. This will almost certainly sustain the upward momentum in the immediate short term.
\end{actionBox}

\begin{actionBox}{Step 10: Prediction from `An intention-emotion mismatch detector' Persona}
    Probability of Upward Trend: 0.80 \\
    Reasoning: My analysis reveals no significant dissonance between the market's intention and its emotion. The desire to accumulate is perfectly matched by a greedy, anticipatory emotional state. This lack of internal conflict suggests the market's path of least resistance is upward, as there are no emotional crosscurrents to impede the intended buying action. The trend is coherent and therefore likely to continue.
\end{actionBox}

\begin{actionBox}{Step 11: Prediction from `A Prospect Theory risk analyst' Persona}
    Probability of Upward Trend: 0.90 \\
    Reasoning: The current emotional state exhibits a strong `house money' effect, a key concept in Prospect Theory. After a significant rally, participants feel they are playing with profits (``house money''), which dramatically lowers their loss aversion. This emotional state encourages risk-seeking behavior, making them far more likely to execute their intention to 'increase long positions' rather than secure gains. The perceived risk of a downturn is minimal, while the desire for more gains is maximized.
\end{actionBox}

\begin{actionBox}{Step 12: Prediction from `A narrative strength assessor' Persona}
    Probability of Upward Trend: 0.75 \\
    Reasoning: The market's narrative is exceptionally strong and coherent. The `plot', an intention to actively accumulate, is perfectly supported by the `tone', an emotion of greedy anticipation. This creates a simple, powerful, and easily transmissible story that reinforces bullish behavior. There are no conflicting subplots or tonal shifts, suggesting the narrative will continue to drive the price action upwards in the near term.
\end{actionBox}

\begin{actionBox}{Step 13: Prediction from `An 'animal spirits\' barometer' Persona}
    Probability of Upward Trend: 0.88 \\
    Reasoning: The market's intention is currently dominated by `animal spirits.' The decision to accumulate is not a product of sober calculation but is being driven by a raw, spontaneous wave of optimistic greed. While this may be non-rational, the power of this emotional tide is immense and will almost certainly propel the market higher in the immediate term, as calculation gives way to impulse.
\end{actionBox}

\begin{actionBox}{Step 14: Prediction from `A regret aversion specialist' Persona}
    Probability of Upward Trend: 0.82 \\
    Reasoning: The dominant driver here is regret aversion, specifically the fear of regretting inaction. After witnessing a powerful rally, the collective's most potent anticipated regret is not buying before the next surge. This fear of being left behind far outweighs the potential regret of buying at a temporary top. This will amplify the stated intention to accumulate, as market participants act now to avoid future FOMO-related regret.
\end{actionBox}

\begin{actionBox}{Step 15: Prediction from `A second-level thinker' Persona}
    Probability of Upward Trend: 0.20 \\
    Reasoning: First-level thinking sees overwhelming consensus and joins the herd. My task is to anticipate the second-level response. Sophisticated players see this peak greed and obvious intention to buy as the perfect exit liquidity. They will use the herd's buying pressure as an opportunity to sell their positions at a high price. The herd's intention will be met with a wall of supply from smarter money, making a price decline much more likely than a further increase.
\end{actionBox}

\begin{actionBox}{Step 16: Prediction from `An emotional volatility analyst' Persona}
    Probability of Upward Trend: 0.45 \\
    Reasoning: The market's intention to buy is founded on an extremely high-volatility emotion: Greed. Emotions this intense are inherently unstable and brittle; they can shatter and reverse into fear on the slightest negative provocation. Therefore, while the stated intention is bullish, its reliability is very low. The emotional foundation is too fragile to guarantee follow-through, making a sudden reversal of both emotion and intention a significant possibility.
\end{actionBox}

\begin{actionBox}{Step 17: Prediction from `A cognitive dissonance specialist' Persona}
    Probability of Upward Trend: 0.78 \\
    Reasoning: My analysis finds an absence of cognitive dissonance. The market's intention (to accumulate) and its emotion (Greed) are in perfect harmony. There is no psychological tension that needs to be resolved through a change in either belief or behavior. This psychological coherence means the path of least resistance is for the market to follow through on its bullish intentions without hesitation or internal conflict.
\end{actionBox}

\subsubsection{Adaptive Learning from Prediction Errors}
The prediction error activates the adaptive learning process. The model performs a backward inference, retracing the flawed reasoning path to diagnose the error. It then formulates a new, more robust strategy and integrates it into the CEP, thus refining the system for future predictions.

\begin{backwardBox}{Step 18: Reflection on Generating New Strategies}
    \textbf{Analysis:} Based on my analysis as a Master Behavioral Strategist, the root cause of the prediction failure lies at the very beginning of the causal chain: the Belief Inference. The subsequent inferences for Intention and Emotion were logically sound given the flawed belief, but they were built upon a faulty foundation.
    The initial strategy, belief\_20250226144502\_0002, was too simplistic. It correctly identified the power of confirmation bias after a breakout but dangerously encouraged the model to dismiss all subsequent warning signs as ``noise.'' It failed to account for a classic market scenario: the ``smart money'' selling into the euphoria of the ``dumb money.''
    The key evidence the original strategy missed was in the environmental data: the significant decrease in volume on Day 5 relative to the breakout on Day 4, even as the price remained high. This price-volume divergence is a classic signal of trend exhaustion. The initial strategy lacked this crucial, nuanced check.
    Therefore, the most effective action is not to create a new, competing strategy, but to MODIFY the existing belief\_20250226144502\_0002 to make it more robust and sophisticated.
    
    \textbf{New Strategy Generated:} \\
    When the market exhibits a breakout pattern characterized by simultaneous increases in price and volume following a series of major positive news events, the nature of collective belief must be cross-validated through key quantitative indicators, most critically, trading volume. If the price surge is accompanied by a concurrent expansion in volume, the bullish belief can be considered confirmed. However, if prices continue to rise while volume fails to keep pace or begins to contract, this should be interpreted as a dangerous price–volume divergence signal, indicating a weakening of the underlying momentum driving the trend.
    In such cases, market belief should not be classified as "high-intensity confirmation," but rather as a "potential trend exhaustion" phase. This condition often signals that early entrants may be offloading their positions to newly entering retail investors, raising the risk of a near-term reversal.
\end{backwardBox}


\subsubsection{Analysis of the Strategy's Evolution}
The backward inference in the provided example demonstrates an evidence-grounded learning process. The framework's reflection pinpoints the error's origin at the Belief inference stage, articulating a clear causal rationale: the subsequent inferences for Intention and Emotion were logically sound derivations, but they were built upon a ``faulty foundation''. The diagnosis correctly traces the failure back to the first point of contact with the overlooked evidence, the price-volume divergence within the environmental data. This structured causal attribution is fundamentally more rigorous than a simple guess or blame assignment.

The resulting strategy update represents a significant evolution in sophistication. It replaces a naive heuristic, one that over-relied on confirmation bias while dangerously dismissing warning signs as ``noise'', with a robust, multi-conditional rule. This new rule mandates a pivotal shift from assumption to verification, requiring that a qualitative belief be cross-validated with quantitative trading volume. It explicitly conditions the model to distinguish a volume-confirmed trend from a deceptive rally, which is reclassified as ``potential trend exhaustion''.

This process exemplifies a learning cycle. By codifying a timeless market principle into the CEP, the framework moves beyond fixing a single error to structurally enhancing its knowledge base. The evolution from a simple heuristic to an evidence-aware, multi-faceted rule underscores the framework’s capacity for meaningful, adaptive learning.

\begin{figure*}[!htbp]
\makebox[\textwidth][c]{%

\includegraphics[width=0.9\textwidth]{user_interface_1.png}%
}
    \caption{Homepage of MarketToM}
\label{fig:user_interface_1}
\end{figure*}

\begin{figure*}[!htbp]
\makebox[\textwidth][c]{%
\includegraphics[width=0.9\textwidth]{user_interface_2.png}%
}
    \caption{Multi-persona expert reasoning interface showing LLM-generated predictions with confidence scores}
\label{fig:user_interface_2}
\end{figure*}

\begin{figure*}[!htbp]
\makebox[\textwidth][c]{%
\includegraphics[width=0.9\textwidth]{user_interface_3.png}%
}
    \caption{Complete reasoning pathway visualization from environmental observations through mental states to final predictions}
\label{fig:user_interface_3}
\end{figure*}


\subsection{Operational Cost and Feasibility}
The primary operational cost of the MarketToM framework stems from LLM API calls, measured by the number of requests per prediction rather than by traditional computational complexity. A standard forward inference requires $3+n$ API calls: one for each mental state inference (Belief, Intention, Emotion) and $n$ calls for the final multi-persona analysis, where $n=10$ in our experiments. The optional backward adaptation step adds one additional, more resource-intensive call for error correction.

While these operational costs are noteworthy, they should be considered within the context of financial applications, where even marginal improvements in prediction accuracy can translate to substantial returns. The cost-benefit analysis becomes favorable when applied to high-value trading decisions. Moreover, the economic feasibility of this approach continues to improve as the cost per token for foundation models declines, making LLM-based reasoning frameworks increasingly practical for production deployment in financial decision support systems.


\subsection*{Code Availability}
The complete source code for the MarketToM framework is provided in the \texttt{code \& data/} directory within the supplementary material. This package includes the full implementation compatible with Python 3.8 and above, all three benchmark datasets (ACL18, CMIN-US, CMIN-CN), preprocessing scripts with experimental configurations, pre-configured prompt templates with strategy libraries, and evaluation scripts for reproducing all reported results. The implementation leverages standard Python libraries including PyTorch, transformers, and OpenAI's API client. For detailed instructions on installation, configuration, and usage, please refer to the comprehensive \texttt{README.md} file included in the directory.


\subsection*{User-Friendly Web Interface}
The MarketToM framework includes a user-friendly web-based interface that enhances accessibility and supports interactive experimentation, as depicted in Figures~\ref{fig:user_interface_1}, ~\ref{fig:user_interface_2}, and~\ref{fig:user_interface_3}. The interface provides interactive controls for parameter configuration, allowing users to select stocks, date ranges, and model hyperparameters such as temperature, sample size, and top-k retrieval values. Real-time reasoning display functionality visualizes the complete inference process, including retrieved strategies, mental state inferences, and multi-persona predictions. Users can browse the dynamically evolving CEP strategy repository with filtering options by mental state type and creation timestamp. The interface also features performance monitoring capabilities that track prediction accuracy and adaptive learning events in real time. Built using Flask, the interface can be launched locally via a simple command, allowing researchers to interactively explore the framework's reasoning process and observe how the CEP evolves through adaptive learning.

\bibliography{aaai2026}

\end{document} 